\subsection{Phase 2: Practical Analyzability}\label{sec:method:source}

The second phase measures the practical boundary of bytecode analysis.  The unit is a \texttt{.pyc} file, not a package: for each bytecode-containing artifact, the analyzer extracts every \texttt{.pyc} entry, resolves its CPython tag, and tests whether selected tools can load, disassemble, or decompile it.  The outcome is the highest level reached: not loadable, loadable only, disassemblable, partially decompilable, or fully decompilable.

The phase is version-aware by design.  Python bytecode is an implementation detail that can change across CPython releases, and the standard disassembler documents this version sensitivity~\cite{python-dis}.  Therefore, a file is included in tool analysis only when the scan identifies a supported modern CPython tag and a matching interpreter is prepared.  In the current scope, this means CPython~3.8 and later; CPython~3.7 and earlier, PyPy tags, unknown tags, and noncanonical tags without an interpreter mapping remain in the raw scan inventory but are excluded from tool-failure counts.  This prevents missing interpreters or unsupported runtimes from being misreported as tool limitations.

The environment-preparation step reads the scanner's version summary, uses the interpreter field to identify canonical CPython runtimes, locates or installs matching interpreters when the host supports it, and prepares per-version virtual environments for tools that depend on bytecode format.  Each analysis step records both status and failure reason, including marshal failure, disassembly failure, unsupported bytecode version, external-tool timeout, nonzero exit, and unavailable decompiler output.  These failure reasons are necessary for interpretation: an environment miss, a loader failure, and a decompiler failure imply different limits.

Tool selection follows five criteria.  \textbf{C1: Availability} requires public, locally runnable tools so the study does not depend on closed services.  \textbf{C2: Python~3 support} keeps the baseline aligned with the modern CPython bytecode observed in the scan.  \textbf{C3: Practical use} favors tools that analysts or developers are likely to encounter through documentation, packaging, or community references.  \textbf{C4: Maintenance or stability} avoids failures caused mainly by obsolete or unreproducible tooling.  \textbf{C5: Method diversity} ensures that the selected set covers loading, disassembly, and decompilation rather than one tool family.  These criteria define a reproducible baseline for this study; they do not rank all Python bytecode tools globally.

\begin{table}[t]
  \centering
  \caption{Selected bytecode-analysis tools and selection criteria.}
  \label{tab:tool-selection}
  \begin{tabular}{@{}llp{0.25\linewidth}p{0.35\linewidth}@{}}
    \toprule
    Tool & Criteria & Functionality & Role \\
    \midrule
    CPython \texttt{marshal}~\cite{python-marshal} & C1,C2,C4,C5 & Loads code objects from \texttt{.pyc} payloads. & Reference code-object loading. \\
    CPython \texttt{dis}~\cite{python-dis} & C1,C2,C4,C5 & Disassembles bytecode instructions. & Standard bytecode disassembly. \\
    \texttt{uncompyle6}~\cite{uncompyle6} & C1--C5 & Decompiles Python bytecode to source. & Established decompilation baseline. \\
    \texttt{decompyle3}~\cite{decompyle3} & C1--C5 & Python~3 decompilation. & Python-3-focused decompiler baseline. \\
    \texttt{pycdc}~\cite{pycdc} & C1--C5 & Disassembles and decompiles bytecode in C++. & Independent C++ decompiler/disassembler. \\
    PyLingual~\cite{pylingual} & C1--C5 & Modern CPython decompilation. & Recent Python~3 decompiler. \\
    \bottomrule
  \end{tabular}
\end{table}

Table~\ref{tab:tool-selection} applies these criteria to the measurement stack.  CPython \texttt{marshal} and \texttt{dis} provide reference loading and disassembly baselines; the decompilers test whether bytecode can be lifted back toward source-level workflows.  We do not claim this set exhausts all Python bytecode tooling.  It is selected to cover the main local analysis levels available to users while avoiding overclaiming about tools not evaluated.
